\subsubsection{Subproceso 3.B: Confirmación y Registro Transaccional Oficial}

\textbf{Fuentes documentales:}
\begin{itemize}
    \item \textbf{Reglamento de Matrícula (RR-0570-2022):} Norma institucional cuyo Capítulo VIII (arts. 75° a 80°) regula la verificación por sección, las rectificaciones de la primera semana, el Cierre de Matrícula, la oportunidad de los padrones, el reporte a SUNEDU y MINEDU y la devolución del registro a las Facultades; incluye la Disposición Transitoria Única que autoriza la automatización de los trámites de matrícula.
    \item \textbf{Ficha TUPA del procedimiento de Matrícula Regular (PA5400954D, RR-3698-2024):} Documento oficial que establece la aprobación automática del trámite, la atención en dos días hábiles y su carácter gratuito.
    \item \textbf{Ley Universitaria N.\textsuperscript{o} 30220 y Estatuto de la UNI:} Marco legal e institucional general del registro académico universitario, base legal del procedimiento conforme a la ficha TUPA.
\end{itemize}

\textbf{Fuentes tecnológicas:}
\begin{itemize}
    \item \textbf{Intranet DIRCE (\texttt{https://dirce.uni.edu.pe/}):} Portal institucional observado para registrar el flujo de comprobantes del registro, el estado de la matrícula del periodo y los mensajes del sistema durante el cierre.
    \item \textbf{Sistema SIGA:} Plataforma de gestión académica cuya lógica transaccional de asentamiento y cierre fue inferida a partir de la observación del comportamiento del portal y de la documentación institucional disponible.
\end{itemize}

\textbf{Fuentes humanas:}
\begin{itemize}
    \item \textbf{Registro de campo del periodo 2026-2 aportado por un integrante del equipo:} Anotaciones elaboradas a partir de la experiencia directa del proceso en el periodo 2026-2 y de comentarios recabados de estudiantes y de la comunidad de Facultad. Esta fuente respalda el diagnóstico del cierre, en particular la comunicación Facultad--DIRCE al concluir la matrícula y la falta de acceso de la Facultad para realizar modificaciones en el registro.
\end{itemize}
